\subsection*{NT-A18--A19. Reciprocal-prime weighting removes the norm-drift obstruction}

Keep the finite Galois extension \(L/K\), conjugacy class
\(C\), cyclic subgroup \(H=\langle g\rangle\), field
\(F=L^H\), and \(\delta=|C|/|G|\) of NT-A12.
Fix unitary Hecke characters \(\eta_1,\ldots,\eta_r\)
of \(K\), and let \(S\) contain all ramification and
conductor primes. Write
\(\nu_i=\eta_i\circ N_{F/K}\).
For this result the independence condition is only
\[
 \prod_i\eta_i^{m_i}\text{ has infinite order}
              \qquad(0\ne\mathbf m\in\mathbf Z^r).         \tag{A18.1}
\]
Nontrivial unitary norm powers are allowed.
The sufficient analytic input is RH for \(\zeta_L\)
and for every primitive, non-pure-norm member of
\[
 L_F\left(s,\chi\prod_i\nu_i^{m_i}\right),
      \quad \chi\in\widehat H,\quad
                         0\ne\mathbf m\in\mathbf Z^r .
                                                               \tag{A18.2}
\]
Pure norm members require only RH for \(\zeta_F\),
already supplied by RH for \(\zeta_L\). Thus \(M\)
supplies exactly the stated rank-one input, without
general Artin holomorphy. All data in this subsection
are fixed; no uniformity as frequencies approach a
norm relation is asserted.

\paragraph{NT-A18: finite harmonic constants and the exact weighted torus gate.}
Let \(\mathcal P_C\) be the primes outside \(S\) with
Frobenius class \(C\). For \(q_{\mathfrak p}=N_K\mathfrak p\),
put
\[
 W_C(x)=\sum_{\substack{\mathfrak p\in\mathcal P_C\\
                                      q_{\mathfrak p}\le x}}
                       q_{\mathfrak p}^{-1},\qquad
 W_{\mathbf m}(x)=
 \sum_{\substack{\mathfrak p\in\mathcal P_C\\
                                      q_{\mathfrak p}\le x}}
       \frac{\prod_i\eta_i(\mathfrak p)^{m_i}}{q_{\mathfrak p}}.
 \]
There are finite constants \(B_C\in\mathbf R\) and
\(B_{\mathbf m}\in\mathbf C\), for every nonzero
\(\mathbf m\), such that
\[
 \begin{aligned}
 W_C(x)&=\delta\log\log x+B_C
                         +O_{L,K,S}\!\left(\frac{\log(2x)}{\sqrt x}\right),\\
 W_{\mathbf m}(x)&=B_{\mathbf m}
          +O_{\mathbf m,\boldsymbol\eta,L,K,S}\!
              \left(\frac1{\log x}+\frac{\log(2x)}{\sqrt x}\right).
 \end{aligned}                                             \tag{A18.3}
\]
If none of the factors for this \(\mathbf m\) in
(A18.2) is pure norm, its \(1/\log x\) term can be
omitted. The series defining \(B_{\mathbf m}\) is
ordered by prime norm; absolute convergence of its
first powers is not claimed.

The probability measures
\[
 \frac1{W_C(x)}
   \sum_{\substack{\mathfrak p\in\mathcal P_C\\q_{\mathfrak p}\le x}}
        \frac{\delta_{(\eta_1(\mathfrak p),\ldots,\eta_r(\mathfrak p))}}
             {q_{\mathfrak p}}
                                                               \tag{A18.4}
\]
therefore tend weakly to Haar measure on \((S^1)^r\).
Every fixed product of arcs has limiting mass equal
to its product length, with arbitrary endpoints.
For each fixed nonzero Fourier frequency the error
is \(O_{\mathbf m}(1/\log\log x)\).
The same conclusions, with possibly changed finite
constants, hold if only absolute-degree-one primes
are counted.

Conversely a finite-order nonzero monomial prevents
full Haar convergence for these weights. Thus, under
\(M\), (A18.1) is the exact condition for full
reciprocal-prime-weighted Haar distribution. It is
strictly weaker than the ordinary-prime condition
in NT-A17.

There is also a convergent Euler-product formulation.
For every nonzero \(\mathbf m\), the products
\[
 \prod_{\substack{\mathfrak p\in\mathcal P_C\\q_{\mathfrak p}\le x}}
 \left(1-\frac{z\prod_i\eta_i(\mathfrak p)^{m_i}}
                    {q_{\mathfrak p}}\right)^{-1}
                                                               \tag{A18.5}
\]
converge locally uniformly for \(|z|<2\) to a
nonvanishing holomorphic function. In particular
their values at \(z=1\) have finite nonzero limits.

\begin{proof}
First consider a single primitive unitary Hecke
character \(\lambda\) of \(F\), with all conductor
primes removed. Additional finitely many primes
may also be removed. If \(\lambda\) is not pure norm,
NT-A08 under its RH gives, for fixed data,
\[
 A_\lambda(u):=\sum_{N_F\mathfrak q\le u}\lambda(\mathfrak q)
                        =O_\lambda(\sqrt u\log(2u)).
 \]
Abel summation, including any primes of norm two,
gives
\[
 \sum_{N_F\mathfrak q\le x}
            \frac{\lambda(\mathfrak q)}{N_F\mathfrak q}
   =\frac{A_\lambda(x)}x+
                         \int_2^x\frac{A_\lambda(u)}{u^2}\,du .
 \]
The integral over \([2,\infty)\) is absolutely
convergent. Its tail and the endpoint term are
\(O_\lambda(\log(2x)/\sqrt x)\). This proves the
claimed finite limit in the non-pure case.

Now let the ideal values of a pure norm character
be \((N_F\mathfrak q)^{i\tau}\), with \(\tau\ne0\).
Under RH for \(\zeta_F\), the prime count after the
same finite deletions is
\(\operatorname{Li}_2(u)+E(u)\), where
\(E(u)=O_{F,S}(\sqrt u\log(2u))\).
Stieltjes integration by parts with
\(u^{i\tau-1}\) gives
\[
 \begin{split}
 \sum_{N_F\mathfrak q\le x}(N_F\mathfrak q)^{i\tau-1}
  ={}&\int_2^x\frac{u^{i\tau-1}}{\log u}\,du
       +x^{i\tau-1}E(x)\\
    &-(i\tau-1)\int_2^x E(u)u^{i\tau-2}\,du .
 \end{split}                                                   \tag{A18.6}
\]
The \(E\)-integral converges absolutely, and its
tail together with the endpoint is
\[
        O_{F,S}\left((1+|\tau|)
                               \frac{\log(2x)}{\sqrt x}\right).
 \]
The main integral becomes
\(\int_{\log2}^{\log x}e^{i\tau y}\,dy/y\).
It converges at infinity by integration by parts,
with the explicit tail bound
\[
       \left|\int_{\log x}^\infty
                      \frac{e^{i\tau y}}y\,dy\right|
                         \le\frac2{|\tau|\log x}.         \tag{A18.7}
\]
These computations prove actual convergence to a
finite constant in the pure norm case too.
There is no application of the shifted-pole
realification argument in this case.

Condition (A18.1) survives norm pullback:
if a pulled-back monomial had finite order \(b\),
composition with the diagonal idele map would
make its original monomial have order dividing
\(b[F:K]\). Thus, for nonzero \(\mathbf m\),
\(\chi\prod_i\nu_i^{m_i}\) is never principal.
If it is pure norm, its parameter \(\tau\) is
therefore nonzero. RH for \(\zeta_L\) supplies
RH for \(\zeta_F\) by the cyclic factorization
\(\zeta_L=\prod_{\chi\in\widehat H}L_F(s,\chi)\):
the nontrivial finite characters are entire, so
a nontrivial zero of any factor cannot be
cancelled in the open critical strip.

Apply the cyclic class-function identity (A12.9)
to the weighted prime sum. The part over primes
of \(F\) with relative residue degree one is
exactly \(W_{\mathbf m}(x)\), with coefficient
\(\delta\overline{\chi(g)}\) in the \(\chi\)-sum.
Each full prime sum over \(F\) has a finite limit
by (A18.6)--(A18.7) or the preceding non-pure
argument. The additional primes of relative
residue degree at least two give an absolutely
convergent correction, uniformly in their unit
character values.

Here is a quantitative check of that correction.
Every such prime has absolute residue degree
at least two. For \(n_F=[F:\mathbf Q]\),
\[
 \begin{split}
 \sum_{\substack{\mathfrak q:\ f_{\mathfrak q}\ge2\\
                                   N_F\mathfrak q>x}}
                \frac1{N_F\mathfrak q}
 &\le n_F\sum_{j=2}^{n_F}\sum_{p>x^{1/j}}p^{-j}\\
 &\le n_F\sum_{j=2}^{n_F}
       \left(x^{-1}+\frac{x^{-1+1/j}}{j-1}\right)
                       \ll_F x^{-1/2}.
 \end{split}                                                   \tag{A18.8}
\]
The inner bound follows by comparison with the
integral of \(t^{-j}\), and the rational primes
may be enlarged to all integers. The empty
sum for \(n_F=1\) causes no exception.
Combining these finitely many \(\chi\)-sums and
corrections proves the second line of (A18.3).

For the zeroth frequency, NT-C02 and finite
deletions give
\[
 \pi_C(u)=\delta\operatorname{Li}_2(u)+R_C(u),
          \qquad R_C(u)=O_{L,K,S}(\sqrt u\log(2u)).
 \]
The complete reciprocal-sum calculation of
NT-C08 applies verbatim: define
\[
 B_C=-\delta\log\log2+
                             \int_2^\infty\frac{R_C(u)}{u^2}\,du .
 \]
Then the error is exactly
\(R_C(x)/x-\int_x^\infty R_C(u)du/u^2\).
This proves the first line of (A18.3), and
\(W_C(x)\to\infty\).
Deleting absolute-degree-at-least-two primes
of \(K\) changes each Fourier sum by a convergent
series with the tail (A18.8), now over \(K\).
This proves every assertion about the degree-one
version.

Divide each nonzero Fourier sum by \(W_C(x)\).
Its limit is zero, with the stated fixed-frequency
rate. For completeness, trigonometric polynomials
are uniformly dense in continuous functions on
the torus, for instance by product Fej\'er means.
For a polynomial the limiting integral is Haar
integration by the just-proved Fourier limits.
For a continuous function approximate uniformly;
both the measures in (A18.4) and Haar measure
have mass one, so the approximation error is
at most twice the uniform error. This proves
weak convergence. Continuous upper and lower
approximations to a fixed box, with arbitrarily
small Haar integral gap at its faces, prove the
box assertion and every endpoint convention.

If instead a nonzero monomial is finite order
dividing \(b\), its \(b\)-th power is identically
one outside the conductor set. The corresponding
nonzero Fourier coefficient of (A18.4) is exactly
one whenever the denominator is positive.
Its Haar integral is zero, which proves necessity.

Finally set \(u_{\mathfrak p}=\prod_i\eta_i(\mathfrak p)^{m_i}\).
For \(|z|<2\), use the power-series logarithm of
each factor in (A18.5). Its first-degree sum is
\(zW_{\mathbf m}(x)\), converging locally uniformly.
For each fixed \(R<2\), all higher-degree sums
are uniformly absolutely convergent on \(|z|\le R\),
since
\[
 \sum_{\mathfrak p\in\mathcal P_C}\sum_{j\ge2}
       \frac{R^j}{j q_{\mathfrak p}^{\,j}}
 \le \frac{R^2}{2(1-R/2)}
                 \sum_{\mathfrak p\in\mathcal P_C}q_{\mathfrak p}^{-2}
                                                    <\infty .
 \]
The resulting logarithm is holomorphic. Its
exponential is the asserted locally uniform
Euler-product limit and never vanishes.
\end{proof}

\paragraph{NT-A19: a Dirichlet-GRH-only progression example.}
Fix primitive Dirichlet characters \(\chi_i\) of
conductors \(f_i\), and real numbers \(t_1,\ldots,t_r\)
which are linearly independent over \(\mathbf Q\).
For primes \(p\nmid\prod_i f_i\), put
\[
                    z_i(p)=\chi_i(p)p^{it_i}.
 \]
Fix a reduced arithmetic progression \(p\equiv a\pmod q\),
and put \(Q_0=\operatorname{lcm}(q,f_1,\ldots,f_r)\).
RH for the primitive Dirichlet characters with
conductor dividing \(Q_0\), including \(\zeta\) for
conductor one, suffices for the following assertions.
This is a finite subfamily of the user-adopted input
\(A_D\), and uses no additional \(M\) input.

The \(1/p\)-weighted points \((z_1(p),\ldots,z_r(p))\)
in that progression tend to full Haar measure on
\((S^1)^r\). Every nonzero mixed prime harmonic sum
converges to a finite constant. In contrast, the
ordinary equal-weight prime points fail to be
Haar equidistributed.

\begin{proof}
Partition the progression into its finitely many
reduced residue classes \(b\bmod Q_0\) satisfying
\(b\equiv a\bmod q\), omitting the finite conductor
primes. On each such class, a Fourier frequency
\(\mathbf m\ne0\) has value
\[
       \prod_i z_i(p)^{m_i}
        =\left(\prod_i\chi_i(b)^{m_i}\right)p^{i\tau_{\mathbf m}},
       \qquad \tau_{\mathbf m}=\sum_i m_i t_i\ne0.
 \]
The progression estimate NT-P12, with the
moduli-one-and-two convention there, gives
\[
 \pi(u;Q_0,b)
   =\frac{\operatorname{Li}_2(u)}{\varphi(Q_0)}
                              +O_{Q_0}(\sqrt u\log(2u)).
 \]
Apply exactly the Stieltjes calculation (A18.6),
with the main integral multiplied by
\(1/\varphi(Q_0)\). Equation (A18.7) proves
convergence of each Fourier harmonic sum on
each class, hence on their finite union.
The zeroth sum is
\(\varphi(q)^{-1}\log\log x+O_{q,\boldsymbol f}(1)\).
The same Fourier approximation argument therefore
proves full weighted Haar distribution.

For the ordinary weighting, choose any coordinate
\(i\) and an integer \(b_i\ge1\) with
\(\chi_i^{b_i}=1\). Its nonzero Fourier frequency
\(b_i\mathbf e_i\) has value \(p^{ib_i t_i}\).
The ordinary Abel calculation in NT-A17, now
using NT-P12 in the fixed progression, gives
\[
 \frac{\sum_{p\le x,\ p\equiv a\bmod q}p^{ib_i t_i}}
               {\pi(x;q,a)}
       =\frac{x^{ib_i t_i}}{1+ib_i t_i}+o(1).
 \]
Omitting conductor primes does not affect the
limit. Its modulus tends to the positive number
\((1+b_i^2t_i^2)^{-1/2}\), whereas the Haar
Fourier integral is zero. This proves the
ordinary-weight obstruction.
\end{proof}

\paragraph{Geometric and analytic meaning.}
The measures are reciprocal-norm measures on
closed points in a specified finite \'{e}tale
Frobenius class, with values in a compact torus
of rank-one Hecke coordinates. Their radial
norm drift is averaged by the logarithmic
measure \(dy/y\); a finite character relation
cannot be averaged away.
For general unitary infinity types this is an
idele-class formulation, not a claimed algebraic
motive realization. No zero-independence or
joint-zero-statistics hypothesis is used.
The constants in (A18.7) can grow when a nonzero
frequency approaches a norm relation, so no
uniform shrinking-box result is inferred here.
Full proofs: \texttt{reciprocal\_prime\_tori.tex},
NT-A18--A19; run \texttt{GRH\_NIGHT/number\_theory}.
Exact classical inputs are NT-A08 (S-NT-IK),
NT-C02 (S-NT-GM), the cyclic induction already
proved in NT-A12, and NT-P12 for the \(A_D\)
example. These are newly proved formulations
in this run, with no literature-priority claim.
